%%% 答案册封底 %%%
{
	\thispagestyle{empty}
	\begin{tikzpicture}[remember picture, overlay]
		% === 背景：暖白底 ===
		\fill[coverlight]
			(current page.south west) rectangle (current page.north east);

		% === 封面镜像：底部橙色横条 ===
		\fill[coverorange]
			(current page.south west) rectangle
			([yshift=2.5cm]current page.south east);
		% 横条内系列名（白色粗体）
		\node[anchor=east, text=white, font=\bfseries\large]
			at ([xshift=-2.8cm, yshift=1.25cm]current page.south east)
			{高中数学之旅};

		% === 封面镜像：右侧橙色竖条 ===
		\fill[coverorange]
			(current page.north east) rectangle
			([xshift=-2cm, yshift=-1.5cm]current page.south east);
		% 竖条顶部细横线
		\draw[coverorange, line width=0.8pt]
			([xshift=-2cm, yshift=-1.5cm]current page.north east) --
			([yshift=-1.5cm]current page.north west);

		% === 左侧内容区 ===
		% 书名（深炭灰色）
		\node[anchor=east, text=coverdark, font=\fontsize{28}{34}\selectfont\bfseries,
			text width=12cm, align=flush right]
			at ([xshift=-2.8cm, yshift=3.5cm]current page.east)
			{高中数学之旅（参考答案）};

		% 英文书名（深琥珀色斜体）
		\node[anchor=east, text=coveramber, font=\normalsize\itshape,
			text width=12cm, align=flush right]
			at ([xshift=-2.8cm, yshift=2cm]current page.east)
			{An Excursion through Secondary School Mathematics: Solutions};

		% 分隔线
		\draw[coverorange, line width=0.4pt]
			([xshift=1.5cm, yshift=1cm]current page.west) --
			([xshift=-2.8cm, yshift=1cm]current page.east);

		% 简介（深炭灰，左对齐）
		\node[anchor=north west, text=coverdark, font=\small,
			text width=12.5cm, align=justify]
			at ([xshift=1.5cm, yshift=0.6cm]current page.west)
			{本书为《高中数学之旅》配套参考答案册，与习题集一一对应，
			提供每道选择题、填空题和解答题的完整解析过程。
			答案讲解注重思路拆解与逻辑链条的完整呈现，
			便于读者在练习后自查订正、加深理解。};

		% 作者（底部横条上方）
		\node[anchor=south east, text=coverdark, font=\small]
			at ([xshift=-2.8cm, yshift=2.8cm]current page.south east)
			{知乎@余命数 \enspace\textbar\enspace Midnight Forever};

		% 版本信息（底部横条上方）
		\node[anchor=south east, text=coveramber, font=\small]
			at ([xshift=-2.8cm, yshift=1.6cm]current page.south east)
			{\today};
	\end{tikzpicture}
}
